\documentclass[aspectratio=169]{beamer}

% ======================
% Packages
% ======================
\usepackage[utf8]{inputenc}
\usepackage[T1]{fontenc}
\usepackage[brazil]{babel}
\usepackage{amsmath, amssymb}
\usepackage{csquotes}
\usepackage{booktabs}
\usepackage{colortbl}
\usepackage{array}

% ======================
% Colors
% ======================
\definecolor{mblack}{RGB}{20, 20, 20}
%\definecolor{mblue}{RGB}{113, 162, 182}
%\definecolor{morange}{RGB}{152, 38, 73}
\definecolor{mblue}{RGB}{152, 38, 73}
\definecolor{morange}{RGB}{113, 162, 182}
\definecolor{mwhite}{RGB}{255, 255, 255}
\definecolor{mcc}{RGB}{30, 100, 180}
\definecolor{mia}{RGB}{20, 140, 80}
\definecolor{mcclight}{RGB}{220, 235, 255}
\definecolor{mialight}{RGB}{210, 245, 225}
% Categorias de disciplinas
\definecolor{catmat}{RGB}{208, 228, 248}
\definecolor{catprog}{RGB}{254, 236, 200}
\definecolor{catia}{RGB}{200, 235, 210}
\definecolor{catproj}{RGB}{235, 215, 245}
\definecolor{catest}{RGB}{255, 245, 200}
\definecolor{cathum}{RGB}{250, 215, 225}
\definecolor{catsw}{RGB}{210, 235, 255}
\definecolor{cattcc}{RGB}{230, 220, 250}
\definecolor{catext}{RGB}{210, 245, 248}
\definecolor{catopt}{RGB}{230, 232, 234}

% ======================
% Theme & Style
% ======================
\usetheme{default}
\usecolortheme{default}
\setbeamercolor{normal text}{fg=mblack, bg=white}
\setbeamercolor{frametitle}{fg=white, bg=mblue}
\setbeamercolor{title}{fg=white}
\setbeamercolor{author}{fg=white}
\setbeamercolor{date}{fg=white}
\setbeamercolor{itemize item}{fg=morange}
\setbeamercolor{itemize subitem}{fg=mblue}
\setbeamercolor{enumerate item}{fg=morange}
\setbeamercolor{enumerate subitem}{fg=mblue}
\setbeamercolor{block title}{fg=white, bg=mblue}
\setbeamercolor{block body}{fg=mblack, bg=mblue!8}
\setbeamercolor{alerted text}{fg=morange}

\setbeamerfont{frametitle}{series=\bfseries, size=\large}
\setbeamerfont{title}{series=\bfseries, size=\LARGE}
\setbeamerfont{block title}{series=\bfseries}

\setbeamertemplate{navigation symbols}{}

% Footer
\setbeamertemplate{footline}{%
  \leavevmode%
  \hbox{%
    \begin{beamercolorbox}[wd=0.7\paperwidth, ht=2.25ex, dp=1ex, leftskip=0.3cm]{normal text}%
      \tiny\color{mblack!50} DECOM -- Apresentação dos Cursos de CC e IA
    \end{beamercolorbox}%
    \begin{beamercolorbox}[wd=0.3\paperwidth, ht=2.25ex, dp=1ex, rightskip=0.3cm, right]{normal text}%
      \tiny\color{mblack!50}\insertframenumber{} / \inserttotalframenumber
    \end{beamercolorbox}%
  }%
}

% Frametitle com barra laranja
\setbeamertemplate{frametitle}{%
  \nointerlineskip%
  \begin{beamercolorbox}[wd=\paperwidth, ht=0.55cm, dp=0.2cm, leftskip=0.5cm]{frametitle}%
    \usebeamerfont{frametitle}\insertframetitle%
  \end{beamercolorbox}%
  \vspace{-0.1cm}%
  \textcolor{morange}{\rule{\paperwidth}{2pt}}%
}

% Title page
\setbeamercolor{titlebar top}{bg=mwhite}
\setbeamercolor{titlebar main}{bg=mblue}
\setbeamercolor{titlebar bottom}{bg=morange}

\setbeamertemplate{title page}{%
  \vspace{0pt}%
  \begin{beamercolorbox}[wd=\paperwidth, ht=0.55cm, dp=0pt]{titlebar top}%
  \end{beamercolorbox}%
  \begin{beamercolorbox}[wd=\paperwidth, ht=3.8cm, dp=0pt, center]{titlebar main}%
    \vspace{0.6cm}
    {\usebeamerfont{title}\color{white}\inserttitle\par}%
    \vspace{0.35cm}
    {\color{white}\small\insertauthor\par}%
    \vspace{0.15cm}
    {\color{white}\footnotesize\insertdate\par}%
    \vspace{0.25cm}
  \end{beamercolorbox}%
  \begin{beamercolorbox}[wd=\paperwidth, ht=1.3cm, dp=0pt]{titlebar bottom}%
  \end{beamercolorbox}%
}

\setlength{\itemsep}{4pt}

\newcommand{\impt}[1]{\textcolor{morange}{\textbf{#1}}}
\newcommand{\ntc}[1]{\textcolor{mblue}{#1}}
\newcommand{\cc}[1]{\textcolor{mcc}{\textbf{#1}}}
\newcommand{\ia}[1]{\textcolor{mia}{\textbf{#1}}}

% Macro para célula da grade
% Uso: \D{cor}{texto}
\newcommand{\D}[2]{%
  \cellcolor{#1}\fontsize{6.5pt}{7.5pt}\selectfont\bfseries #2}

% ======================
% Document Info
% ======================
\title{Aula 1 --- Apresentação dos Cursos}
\author{Departamento de Computação (DECOM) -- UFOP}
\date{\today}

% ======================
% Document
% ======================
\begin{document}

% --- Capa ---
\begin{frame}[plain]
  \titlepage
\end{frame}

% --- Sumário ---
\begin{frame}{Sumário}
  \tableofcontents
\end{frame}

% ==============================================================
\section{Contexto Institucional}
% ==============================================================

\begin{frame}{UFOP e o DECOM}
  \begin{itemize}
    \item Ambos os cursos são oferecidos pela \textbf{Universidade Federal de Ouro Preto (UFOP)}, campus Morro do Cruzeiro, Ouro Preto -- MG.
    \item Vinculados ao \textbf{Instituto de Ciências Exatas e Biológicas (ICEB)}, que reúne departamentos das áreas de Exatas, Biológicas e Computação.
    \item O \impt{Departamento de Computação (DECOM)} é o departamento central de ambos os cursos:
    \begin{itemize}
      \item Mais de 30 professores doutores em dedicação exclusiva
      \item Áreas: algoritmos, IA, engenharia de software, redes, otimização, computação gráfica...
    \end{itemize}
  \end{itemize}
  \begin{block}{Canal oficial do aluno}
    O \textbf{Colegiado do Curso} é responsável pela parte educacional e pedagógica dos cursos.
    \centering
    cocic@ufop.edu.br, cobia@ufop.edu.br
  \end{block}
\end{frame}

% ==============================================================
\section{Os Dois Cursos em Perspectiva Comparada}
% ==============================================================
\centering
\begin{frame}{Os Dois Cursos: Números}
  \vspace{0.2cm}
  \renewcommand{\arraystretch}{1.3}
  \begin{tabular}{lcc}
    \toprule
    \textbf{Característica} & \cc{CC} & \ia{IA} \\
    \midrule
    Início              & 1986          & 2025 \\
    Vagas / semestre    & 40            & 40 \\
    Integralização mín. & 8 semestres   & 8 semestres \\
    Integralização máx. & 12 semestres  & 12 semestres \\
    Carga total         & 3.300 h       & 2.820 h \\
    Obrigatórias        & 2.670 h       & 2.100 h \\
    Eletivas (mín.)     & 360 h         & 480 h \\
    Extensionistas      & 330 h         & 360 h \\
    AACC                & 240 h         & 240 h \\
    TCC                 & Monografia I+II & Memorial\\
    \bottomrule
  \end{tabular}
  \vspace{0.2cm}

  \ntc{IA tem carga eletiva maior: formação especializada com maior flexibilidade de aprofundamento.}
\end{frame}

% ==============================================================
\section{Bacharelado em Ciência da Computação}
% ==============================================================

\begin{frame}{CC: Missão e Eixos de Formação}
  \begin{block}{Missão}
    Formar profissionais \impt{aptos a especificar, projetar, implementar e avaliar sistemas computacionais}, de forma inovadora e empreendedora, atendendo às necessidades da sociedade e do mercado.
  \end{block}
  \vspace{0.2cm}
  \textbf{7 Eixos de Formação (SBC):}
  \begin{columns}[T]
    \begin{column}{0.48\textwidth}
      \begin{enumerate}\setlength\itemsep{2pt}
        \item Resolução de Problemas
        \item Desenvolvimento de Sistemas
        \item Desenvolvimento de Projetos
        \item Implantação de Sistemas
      \end{enumerate}
    \end{column}
    \begin{column}{0.48\textwidth}
      \begin{enumerate}\setlength\itemsep{2pt}
        \setcounter{enumi}{4}
        \item Gestão de Infraestrutura
        \item Aprendizado Contínuo e Autônomo
        \item Ciência, Tecnologia e Inovação
      \end{enumerate}
    \end{column}
  \end{columns}
\end{frame}

\begin{frame}{CC: Perfil do Egresso}
  Segundo as Diretrizes Curriculares Nacionais (DCNs), espera-se que o egresso de CC:
  \begin{itemize}
    \item Possua \textbf{sólida formação} em CC e Matemática, capaz de construir aplicativos, ferramentas e infraestrutura de software
    \item Tenha \textbf{visão global e interdisciplinar} de sistemas, transcendendo detalhes de implementação
    \item Conheça a estrutura dos sistemas de computação e os processos de construção e análise
    \item Seja capaz de agir \textbf{reflexivamente}: sistemas atingem direta ou indiretamente a sociedade
    \item Reconheça a importância da \textbf{inovação e criatividade} e as perspectivas de negócio
  \end{itemize}
\end{frame}

%%--------------------------------------------

\begin{frame}{CC: Grade Curricular}
  \vspace{0.1cm}
  \centering
  \setlength{\tabcolsep}{2pt}
  \renewcommand{\arraystretch}{1.1}
  \setlength{\arrayrulewidth}{0.9pt}
  \begin{tabular}{|p{1.5cm}|p{1.5cm}|p{1.5cm}|p{1.5cm}|p{1.5cm}|p{1.5cm}|p{1.5cm}|p{1.5cm}|}
  \hline
  \cellcolor{mcc}\color{white}\bfseries\tiny 1º Período &
  \cellcolor{mcc}\color{white}\bfseries\tiny 2º Período &
  \cellcolor{mcc}\color{white}\bfseries\tiny 3º Período &
  \cellcolor{mcc}\color{white}\bfseries\tiny 4º Período &
  \cellcolor{mcc}\color{white}\bfseries\tiny 5º Período &
  \cellcolor{mcc}\color{white}\bfseries\tiny 6º Período &
  \cellcolor{mcc}\color{white}\bfseries\tiny 7º Período &
  \cellcolor{mcc}\color{white}\bfseries\tiny 8º Período \\
  \hline
  \D{catprog}{Intro. Programação} &
  \D{catprog}{Mat. Discreta I} &
  \D{catprog}{Mat. Discreta II} &
  \D{catprog}{Teoria dos Grafos} &
  \D{catprog}{P\&A Algoritmos} &
  \D{catia}{Proc. Digital Imagens} &
  \D{cattcc}{Monografia I} &
  \D{cattcc}{Monografia II} \\
  \hline
  \D{catprog}{Eletrônica p/ Comp.} &
  \D{catprog}{Estrut. de Dados I} &
  \D{catprog}{Estrut. de Dados II} &
  \D{catprog}{Sist. Operacionais} &
  \D{catprog}{Teoria Computação} &
  \D{catprog}{Sist. Distribuídos} &
  \D{catext}{Eventos para CC (120h)} &
  \D{cathum}{Hist. da Filosofia} \\
  \hline
  \D{catprog}{Intro. à CC} &
  \D{catprog}{Org. Computadores} &
  \D{catprog}{POO} &
  \D{catprog}{Eng. de Software I} &
  \D{catprog}{Redes Computadores} &
  \D{catia}{Intelig. Artificial} &
  \D{catext}{Informática e Soc.(90h)} &
  \D{cathum}{Direito Informática} \\
  \hline
  \D{catmat}{Cálculo I} &
  \D{cathum}{IHC} &
  \D{catprog}{Prog. Funcional} &
  \D{catprog}{Cálculo Numérico} &
  \D{catia}{Intro. Aprend. Máquina} &
  \D{cathum}{Metod. Científica} &
  \D{catprog}{Computação Gráfica} &
   \\
  \hline
  \D{catmat}{Geom. Analítica} &
  \D{cathum}{Prát. Leitura/Texto} &
  \D{catprog}{Arq. Computadores} &
  \D{catprog}{Banco de Dados I} &
  \D{catprog}{Eng. de Software II} &
  \D{catprog}{Const. Compiladores I} &
   & \\
  \hline
  &
  \D{catmat}{Cálculo II} &
  \D{catmat}{Álgebra Linear} &
  \D{catmat}{Estat. e Prob.} &
  \D{catprog}{Intro. Otimização} &
  & & \\
  \hline
  \end{tabular}

  \vspace{0.15cm}
  {\tiny
  \colorbox{catmat}{\phantom{x}}~Mat./Est.\quad
  \colorbox{catprog}{\phantom{x}}~CC\quad
  \colorbox{catia}{\phantom{x}}~IA/ML\quad
  \colorbox{cathum}{\phantom{x}}~Humanidades\quad
  \colorbox{cattcc}{\phantom{x}}~TCC\quad
  \colorbox{catext}{\phantom{x}}~Extensão\quad
  }
\end{frame}

% ==============================================================
\section{Bacharelado em Inteligência Artificial}
% ==============================================================

\begin{frame}{IA: Missão e Competências Específicas}
  \begin{block}{Missão}
    Capacitar egressos para \impt{conceber, desenvolver e aplicar sistemas inteligentes} --- alinhados ao Plano Brasileiro de IA (PBIA): \enquote{IA para o Bem de Todos}.
  \end{block}
  \vspace{0.15cm}
  \textbf{5 Competências Específicas do Egresso (CE):}
  \begin{itemize}\setlength\itemsep{3pt}
    \item \textbf{CE-I:} Fundamentos em CC, Matemática e Estatística
    \item \textbf{CE-II:} Paradigmas de IA --- abordagens simbólicas e conexionistas
    \item \textbf{CE-III:} Ciclo de vida dos dados: aquisição, mineração, visualização e ML
    \item \textbf{CE-IV:} Soluções inovadoras com visão de negócio em IA
    \item \textbf{CE-V:} Ética: privacidade, equidade, transparência e vieses algorítmicos
  \end{itemize}
\end{frame}

\begin{frame}{IA: Eixos de Formação}
  \textbf{7 Eixos de Formação específicos para IA (SBC):}
  \begin{columns}[T]
    \begin{column}{0.48\textwidth}
      \begin{enumerate}\setlength\itemsep{3pt}
        \item Fundamentos Matemáticos e Computacionais
        \item Raciocínio e Representação de Conhecimento
        \item Aprendizado de Máquina e Ciência de Dados
        \item Percepção e Atuação (Visão, PLN, Robótica)
      \end{enumerate}
    \end{column}
    \begin{column}{0.48\textwidth}
      \begin{enumerate}\setlength\itemsep{3pt}
        \setcounter{enumi}{4}
        \item Engenharia de Sistemas Inteligentes
        \item Ética, Sociedade e Governança da IA
        \item Desenvolvimento Profissional e Aprendizado Contínuo
      \end{enumerate}
    \end{column}
  \end{columns}
  \pause
  \vspace{0.3cm}
  \begin{alertblock}{Diferencial: Projetos 1 a 4}
    Do 4º ao 7º período, o aluno desenvolve \textbf{4 disciplinas de Projeto} (cada uma com 90h, sendo 90h extensionistas) aplicando os conhecimentos adquiridos a problemas reais.
  \end{alertblock}
\end{frame}

\begin{frame}{IA: Grade Curricular}
  \vspace{0.1cm}
  \centering
  \setlength{\tabcolsep}{2pt}
  \renewcommand{\arraystretch}{1.1}
  \setlength{\arrayrulewidth}{0.9pt}
  \begin{tabular}{|p{1.5cm}|p{1.5cm}|p{1.5cm}|p{1.5cm}|p{1.5cm}|p{1.5cm}|p{1.5cm}|p{1.5cm}|}
  \hline
  \cellcolor{mia}\color{white}\bfseries\tiny 1º Período &
  \cellcolor{mia}\color{white}\bfseries\tiny 2º Período &
  \cellcolor{mia}\color{white}\bfseries\tiny 3º Período &
  \cellcolor{mia}\color{white}\bfseries\tiny 4º Período &
  \cellcolor{mia}\color{white}\bfseries\tiny 5º Período &
  \cellcolor{mia}\color{white}\bfseries\tiny 6º Período &
  \cellcolor{mia}\color{white}\bfseries\tiny 7º Período &
  \cellcolor{mia}\color{white}\bfseries\tiny 8º Período \\
  \hline
  \D{catia}{Intro. à IA} &
  \D{catia}{IA Clássica} &
  \D{catia}{Intro. Ciência Dados} &
  \D{catprog}{Eng. de Software I} &
  \D{catprog}{Teoria dos Grafos} &
  \D{catia}{Proc. Ling. Natural} &
  \D{catmat}{Análise Séries Temp.} &
  \D{cathum}{Seg., Ética e Soc. IA} \\
  \hline
  \D{catprog}{Intro. Programação} &
  \D{catprog}{Estrut. de Dados I} &
  \D{catia}{AM Supervisionado} &
  \D{catia}{Redes Neurais / DL} &
  \D{catia}{Aprend. por Reforço} &
  \D{catia}{Visão Computacional} &
  \D{catproj}{Projeto 4 } &
  \D{cattcc}{Monografia (TCC)} \\
  \hline
  \D{catprog}{Mat. Discreta I} &
  \D{catprog}{Fund. Arq. e Org.} &
  \D{catia}{AM Não Supervis.} &
  \D{catprog}{Intro. Otimização} &
  \D{catprog}{Comp. Paralela/Dist.} &
  \D{cathum}{Metod. Cient. IA} &
   &
   \\
  \hline
  \D{catmat}{Cálculo A} &
  \D{catmat}{Estat. e Prob.} &
  \D{catprog}{POO} &
  \D{catprog}{Intelig. de Negócios} &
  \D{catest}{Planej. Experimentos I} &
  \D{catproj}{Projeto 3 } &
   &
   \\
  \hline
  \D{catmat}{Geom. Anal. e Álg. Lin.} &
  \D{catmat}{Cálculo B} &
  \D{catprog}{P\&A Algoritmos} &
  \D{catproj}{Projeto 1}  &
  \D{catproj}{Projeto 2}  &
    &
    &
   \\
  \hline
  \end{tabular}

  \vspace{0.15cm}
  {\tiny
  \colorbox{catmat}{\phantom{x}}~Mat./Est.\quad
  \colorbox{catprog}{\phantom{x}}~CC\quad
  \colorbox{catia}{\phantom{x}}~IA/ML\quad
  \colorbox{catproj}{\phantom{x}}~Extensão\quad
  \colorbox{cathum}{\phantom{x}}~Humanística\quad
  \colorbox{cattcc}{\phantom{x}}~TCC\quad
  }
\end{frame}

% ==============================================================
\section{Departamentos Envolvidos}
% ==============================================================

\begin{frame}{Departamentos Envolvidos nos Cursos}
  \begin{columns}[T]
    \begin{column}{0.48\textwidth}
      \textcolor{mcc}{\textbf{Colegiado de CC (COCIC)}}
      \begin{itemize}\setlength\itemsep{2pt}
        \item \textbf{DECOM} --- disciplinas obrigatórias e eletivas
        \item \textbf{DEMAT} --- Cálculo I/II, Álgebra Linear, Geom. Anal.
        \item \textbf{DEEST} --- Estatística e Probabilidade
        \item \textbf{DEFIL} --- Intro. à História da Filosofia
        \item \textbf{DEDIR} --- Direito da Informática
        \item \textbf{CEAD} --- Prática de Leitura e Produção de Textos
      \end{itemize}
    \end{column}
    \begin{column}{0.48\textwidth}
      \textcolor{mia}{\textbf{Colegiado de IA}}
      \begin{itemize}\setlength\itemsep{2pt}
        \item \textbf{DECOM} --- disciplinas obrigatórias e eletivas
        \item \textbf{DEMAT} --- Cálculo A/B, Geom. Anal. e Álg. Linear
        \item \textbf{DEEST} --- Estat., Planej. Experimentos, Séries Temporais
      \end{itemize}
    \end{column}
  \end{columns}
  \pause
  \vspace{0.3cm}
  \footnotesize
  \begin{alertblock}{Dica}
    Para aproveitamento de disciplinas, reopção ou qualquer questão acadêmica, o canal correto é o \textbf{Colegiado do seu curso}. Todos os departamentos parceiros têm representação nesses colegiados.
  \end{alertblock}
\end{frame}

% ==============================================================
\section{Como os Cursos se Relacionam}
% ==============================================================

\begin{frame}{CC e IA: Semelhanças e Diferenças}
  \begin{columns}[T]
    \begin{column}{0.48\textwidth}
      \begin{block}{Em comum}
        \begin{itemize}\setlength\itemsep{2pt}
          \item Mesmo DECOM, docentes e laboratórios
          \item Base matemática: cálculo, álgebra, mat.\ discreta
          \item Disciplinas BCC compartilhadas: ED, POO, Grafos, Algoritmos, Redes Neurais...
          \item Empresa Júnior \textbf{Voluta Soluções Digitais}
          \item Acesso ao PPGCC (mestrado e doutorado)
          \item Possibilidade de reopção entre cursos
        \end{itemize}
      \end{block}
    \end{column}
    \begin{column}{0.48\textwidth}
      \cc{CC --- Formação generalista}
      \begin{itemize}\setlength\itemsep{1pt}
        \item Hardware, SO, compiladores, redes, IA, computação gráfica...
        \item 3.300h totais
      \end{itemize}
      \vspace{0.3cm}
      \ia{IA --- Formação especializada}
      \begin{itemize}\setlength\itemsep{1pt}
        \item Toda a grade converge para sistemas inteligentes
        \item 4 disciplinas de Projeto com problemas reais
        \item 480h de eletivas para aprofundamento
      \end{itemize}
    \end{column}
  \end{columns}
\end{frame}

% ==============================================================
\section{Vida Acadêmica}
% ==============================================================

\begin{frame}{Oportunidades Além da Sala de Aula}
  \begin{columns}[T]
    \begin{column}{0.48\textwidth}
      \textbf{Pesquisa e Extensão}
      \begin{itemize}\setlength\itemsep{3pt}
        \item \textbf{Iniciação Científica (IC)} --- bolsas CNPq, FAPEMIG, PIBIC/UFOP
        \item \textbf{Atividades Extensionistas} --- obrigatórias (10\% da carga horária)
        \item \textbf{AACCs} --- eventos, cursos, monitorias, publicações
      \end{itemize}
    \end{column}
    \begin{column}{0.48\textwidth}
      \textbf{Empresa Júnior}
      \begin{itemize}\setlength\itemsep{3pt}
        \item \textbf{Voluta Soluções Digitais} --- fundada em 2016, aberta a alunos de CC
        \item Projetos reais para micro e pequenas empresas
        \item Desenvolve visão empreendedora e gestão
      \end{itemize}
      \vspace{0.2cm}
      \textbf{Pós-Graduação}
      \begin{itemize}
        \item \textbf{PPGCC} --- mestrado e doutorado, aberto a alunos de ambos os cursos
      \end{itemize}
    \end{column}
  \end{columns}
\end{frame}

% ==============================================================
\section{Síntese}
% ==============================================================

\begin{frame}{O Que Você Precisa Saber ao Ingressar}
  \begin{enumerate}\setlength\itemsep{5pt}
    \item \impt{Curso sólido} --- IFES com tradição em pesquisa, extensão e ótimas avaliações no SINAES
    \item \ntc{Múltiplos departamentos} --- professores do DECOM, DEMAT, DEEST e outros ao longo do curso
    \item \textbf{Regime semestral} --- até 6 disciplinas por semestre a partir do 2º período; respeite os pré-requisitos
    \item \textbf{Eletivas e flexibilidade} --- CC: 360h + 60h facultativas; IA: 480h eletivas
    \item \textbf{Pesquisa e extensão} --- bolsas, Voluta, PPGCC à sua disposição
    \item \textbf{Planeje o TCC cedo} --- CC: inicia no 7º período; IA: Projetos 1--4 constroem o TCC progressivamente
  \end{enumerate}
\end{frame}

\end{document}